% TOP_PROBLEM: {"id": 5300016, "title": "Combinatorial boundary of Blaschke-product space", "category_id": 11, "category": {"id": 11, "name": "dynamical_systems", "display_name": "Dynamical Systems", "description": "Problems about long-term behavior of deterministic systems, Hamiltonian dynamics, and periodic orbits.", "slug": "dynamical-systems", "order_index": 11, "created_at": "2026-07-31T15:26:25.671Z"}, "status": "open"}
% FIRST_POSTED: null
% ENTRY_KIND: statement_only
% Imported from: batch04.tex; UnsolvedMath ID 5300016
\documentclass[11pt]{article}
\usepackage{fontspec}
\setmainfont{Latin Modern Roman}
\setsansfont{Latin Modern Sans}
\usepackage{amsmath,amssymb,mathtools}
\usepackage{unicode-math}
\setmathfont{Latin Modern Math}
\usepackage{xurl}
\usepackage[hidelinks]{hyperref}
\newcommand{\sref}[2]{\hyperlink{source:#1:#2}{[S#2]}}
\newcommand{\cataloguescope}[1]{\par\noindent\textbf{Question.} #1\par}
\begin{document}
% UnsolvedMath ID 5300016; source catalogue code AMR-052-0016
\setcounter{section}{5300015}
\section{Combinatorial boundary of Blaschke-product space}\label{5300016}
{\small\sffamily UnsolvedMath ID 5300016 · Dynamical Systems}
\subsection{Definitions and mathematical statement}

\paragraph{Shared notation.}$\mathbb N=\{1,2,\ldots\}$, and $\mathbb Z,\mathbb Q,\mathbb R,\mathbb C$ denote the integers, rationals, reals and complex numbers. Any other conventions are specified locally.


Fix $n>1$. A finite Blaschke product of degree $n$ is a proper holomorphic map of the unit disk $\Delta=\{z\in\mathbb C:|z|<1\}$ to itself of that degree. Let $A,B$ be such maps with $A(0)=B(0)=0$. Choose boundary fixed points $a,b$; there is a unique quasisymmetric conjugacy $h$ of the boundary circles with $h(a)=b$. Glue the disks by $h$ and transport the two dynamics to obtain a rational map $f(A,B)$ of the Riemann sphere, defined up to conformal conjugacy. Its Julia set is a quasicircle and the maps on the two complementary disks are conformally conjugate to $A$ and $B$. A quasicircle is a quasiconformal image of a circle. As in the source, the spaces below represent conformal-conjugacy classes; the boundary fixed-point choices are kept consistent, and additional finite combinatorial markings are suppressed.
Write $\mathcal B(z^n)$ for the component of the degree-$n$ polynomial space containing $z^n$ in which there is an attracting fixed point whose immediate basin contains every finite critical point. Equivalently it is the space of maps $f(z^n,B)$. For fixed $A$, put $\mathcal B(A)=\{f(A,B):B\text{ varies}\}$. These spaces inherit complex structures from the polynomial or rational-map parameter space. Their actual boundaries are taken by closure in those degree-$n$ map spaces, with the same conformal-conjugacy convention. The natural biholomorphism is $F(f(z^n,B))=f(A,B)$.
On the actual boundary of $\mathcal B(A)$ identify $f$ and $g$ when they are quasiconformally conjugate, meaning $u\circ f=g\circ u$ for a quasiconformal homeomorphism $u$ of the sphere. Write $\partial(A)$ for the resulting space with its quotient topology. The source equivalently describes these equivalence classes by connected complex submanifolds of the boundary; the quotient is non-Hausdorff for $n>2$.
\paragraph{Statement recorded in the supplied source.}
Give a combinatorial description of the topological space $\partial(z^n)$. Such a description may use laminations, that is, noncrossing collections of chords or leaves encoding identifications, but the source does not require this particular form of model. \sref{5300016}{2}
\subsection{Short English statement}
Describe the quotient boundary using combinatorial data, possibly laminations.
\subsection{Sources}
\begin{description}
\item[\hypertarget{source:5300016:1}{S1}] ULAM.AI and the UnsolvedMath Contributors, UnsolvedMath: A Curated Collection of Open Mathematics Problems, record 5300016 (AMR-052-0016). CC BY 4.0. \url{https://huggingface.co/datasets/ulamai/UnsolvedMath}.
\item[\hypertarget{source:5300016:2}{S2}] Curt McMullen, Rational maps and Teichmüller space, in Ben Bielefeld and Mikhail Lyubich (editors), Problems in holomorphic dynamics (1992). Problem following the quotient-boundary conjecture, p. 9. \url{https://arxiv.org/pdf/math/9205209}.
\end{description}
\label{end:5300016}
\end{document}
